\section{模曲线与模空间}
\label{sec:modular-curves-moduli}
% ============================================================================

\textbf{为什么需要模曲线？}

\cref{cor:torus-tau} 告诉我们：$Y(1) = \SL_2(\Z) \backslash \HH$
参数化所有椭圆曲线的同构类，通过 $j$-不变量与 $\C$ 同构。
这很好——但有时我们不只想分类椭圆曲线，
还想分类"带有额外数据"的椭圆曲线。

\textbf{典型的"额外数据"是什么？}
以 Wiles 证明 FLT 为例。核心步骤是：把"级别 $N$ 的模形式"与"有理椭圆曲线"对应起来。
这里"级别 $N$"就意味着椭圆曲线上记录了某种 $N$ 阶结构（具体是一个 $N$ 阶循环子群）。
要参数化"带 $N$ 阶子群的椭圆曲线"，$Y(1)$ 就不够用了——它把有无额外结构的曲线混为一谈。

\textbf{解决方案：用更小的对称群。}
如果我们只用 $\Gamma_0(N)$ 的等价关系来识别椭圆曲线
（即只把"$N$ 阶子群相同"的曲线视为等价），
就得到一个更细的分类空间 $Y_0(N) = \Gamma_0(N) \backslash \HH$。
不同的 $N$ 和不同的子群类型（$\Gamma_0, \Gamma_1, \Gamma$）
给出不同"粗细程度"的分类。

这些分类空间——\textbf{模曲线}——不只是集合，
它们还具有代数曲线的结构（将在\cref{ch:riemann-surfaces}中看到这是紧黎曼面），
可以定义在有理数甚至整数上，成为研究椭圆曲线算术的核心工具。

\textbf{谷山-志村猜想}（现在是定理）断言：每条有理椭圆曲线 $E$ 都是某个模曲线 $X_0(N)$ 的"商"，
即存在有理映射 $X_0(N) \to E$。这就是 FLT 证明的骨架。

\begin{definition}[模曲线]
\label{def:modular-curves}
对同余子群 $\Gamma$，定义\textbf{（开）模曲线}
\[
  Y(\Gamma) = \Gamma \backslash \HH.
\]
特别地，$Y_0(N) = \GammaO(N) \backslash \HH$，
$Y_1(N) = \GammaI(N) \backslash \HH$，
$Y(N) = \Gamma(N) \backslash \HH$。

通过添加尖点得到\textbf{（紧化）模曲线}
\[
  X(\Gamma) = \Gamma \backslash \HH^*.
\]
$X(\Gamma)$ 是紧黎曼面（将在\cref{ch:riemann-surfaces}中详细讨论）。
\end{definition}

\begin{theorem}[模诠释]
\label{thm:moduli-interpretation}
模曲线参数化带附加结构的椭圆曲线（在 $\C$ 上）：
\begin{enumerate}
  \item $Y(1) = \SL_2(\Z) \backslash \HH$ 参数化复环面的同构类
    （\cref{cor:torus-tau}）。通过 $j$-不变量，$Y(1) \cong \C$。
  \item $Y_0(N)$ 参数化 $\{(E, C)\}$ 的同构类，其中 $E$ 是椭圆曲线，
    $C \subset E$ 是 $N$ 阶循环子群（等价于核为 $\Z/N\Z$ 的同源
    $E \to E/C$）。
  \item $Y_1(N)$ 参数化 $\{(E, P)\}$ 的同构类，其中 $P \in E$ 是
    $N$ 阶点。
  \item $Y(N)$ 参数化 $\{(E, (P, Q))\}$ 的同构类，其中 $(P, Q)$ 是
    $E[N]$ 的一组有序基。
\end{enumerate}
\end{theorem}

\begin{proof}[解释]
\textbf{(2)}：设 $\tau \in \HH$，对应环面 $E_\tau = \C/\Lambda_\tau$
（$\Lambda_\tau = \Z\tau + \Z$）。取 $C = \langle 1/N \rangle \subset E_\tau$，
即 $C = \frac{1}{N}\Z/\Z$ 是 $N$ 阶循环子群。

何时 $(E_\tau, C_\tau) \cong (E_{\tau'}, C_{\tau'})$？
需要 $\alpha \in \C^*$ 使得 $\alpha\Lambda_\tau = \Lambda_{\tau'}$
且 $\alpha C_\tau = C_{\tau'}$。写出条件：
$\alpha\tau = a\tau' + b$，$\alpha = c\tau' + d$（$ad - bc = 1$），
且 $\alpha \cdot \frac{1}{N} \equiv \frac{j}{N} \pmod{\Lambda_{\tau'}}$
对某 $j$ 与 $N$ 互素。

这等价于 $\frac{c\tau' + d}{N} \in \frac{1}{N}\Z + \Lambda_{\tau'}$，
即 $c \equiv 0 \pmod{N}$。
故等价关系正是 $\GammaO(N)$ 的作用。

\textbf{(3) 和 (4)}：类似分析。
\end{proof}


% ============================================================================
